\chapter*{Preface — Read This First}
\addcontentsline{toc}{chapter}{Preface — Read This First}
\markboth{Preface}{Preface}

\section*{How this manual is organised}
\addcontentsline{toc}{section}{How this manual is organised}

\begin{itemize}
	\item \textbf{Part I — Orientation.} What OLK is, who it is for, and the shape of the system at a glance. Read this even if you only have ten minutes.
	\item \textbf{Part II — Environment \& Workflow.} How to get the project running on your machine — including the local Docker/Supabase stack — and the conventions the project expects you to follow.
	\item \textbf{Part III — The Frontend.} Next.js routing, every page, every shared component, and the styling system.
	\item \textbf{Part IV — The Backend.} Supabase as a concept, then the full database schema, every function and trigger, the row‑level security model, and how server‑side actions and notifications work.
	\item \textbf{Part V — Key Workflows End to End.} The same system traced through its three central journeys: exchanging a book, matching a wishlist, and moderating the community — reading top to bottom instead of file to file.
	\item \textbf{Part VI — Operating in Production.} Deploying to Vercel, the full environment variable reference, a maintenance playbook of “if you need to do X, start here,” a candid list of known issues, and a troubleshooting guide.
	\item \textbf{Appendices.} A compact table reference, a command cheat‑sheet, and a glossary.
\end{itemize}

\section*{Conventions used in this manual}
\addcontentsline{toc}{section}{Conventions used in this manual}

Four recurring boxes appear throughout. They are not decoration — they are load‑bearing:

\begin{olknote}
	Background context or a clarification that helps the surrounding text make sense. Safe to skim on a first read.
\end{olknote}

\begin{olktip}
	A shortcut, a command, or a habit that will save you time. Worth remembering.
\end{olktip}

\begin{olkwarn}
	A place where the obvious thing to do is the wrong thing to do. Read these before you touch the related code.
\end{olkwarn}

\begin{olkgotcha}
	A documented, real, currently‑existing bug, gap, or piece of technical debt in the codebase as of this edition — not hypothetical. Each one names the exact file and line. See also Chapter~22, \emph{Known Issues \& Technical Debt}, which aims to collect all of them in one place.
\end{olkgotcha}

\vspace*{4cm}
\begin{olktask}
	A hands‑on exercise. If you are an intern working through this manual for the first time, do these — reading about a codebase and changing it are different skills, and these are designed to be small enough to finish in under an hour each while touching a real, working part of the system.
\end{olktask}

Code listings are copied verbatim from the repository and captioned with their file path, so you can always locate the original with a simple search. Where a listing has been trimmed for length, an ellipsis comment (\texttt{// ...}) marks the cut.

\section*{A note on honesty and openness}
\addcontentsline{toc}{section}{A note on honesty and openness}

This manual documents OLK as it actually exists today — including the unfinished edges, the occasional inconsistency, and the deliberate trade‑offs made along the way. We believe that a guide describing only an idealised design, while the shipped code tells a different story, actively misleads the next person who trusts it. By being candid about where we stand, we hope to save you time and frustration, and to invite you into a culture of transparency.

From the very beginning, OLK has been conceived as an open‑source project. It is not a closed product delivered from on high; it is a living foundation for experimentation and collective growth. We actively encourage students, developers, and technologists to fork the repository, try out emerging tools, propose novel architectures, and push the platform in directions we have not yet imagined. Whether you are refining a database query, prototyping a new feature, or rethinking the entire UI layer, your contributions are welcome here. This manual is your starting point.

\vfill

\begin{flushright}
	First edition — compiled July 2026 \\
	Owais Saleem
\end{flushright}